\runninghead{Flight without Turning Back}
\properlane{flight-without-turning-back}{Not Coming Down from the Housetop: The Commands of Flight as a Rule of Life}

The first commands of the Gospel are practical and urgent: \latin{tunc qui in
Iudǽa sunt, fúgiant ad montes: et qui in tecto, non descéndat tóllere áliquid
de domo sua: et qui in agro, non revertátur tóllere túnicam suam}, ``Then they
that are in Judea, let them flee to the mountains: And he that is on the
housetop, let him not come down to take any thing out of his house: And he that
is in the field, let him not go back to take his coat'' (Mt 24:16--18). They
are followed by a woe on women with child and nursing and the command to pray
that the flight be not in winter or on the sabbath (vv.~19--20). St John
Chrysostom reads these words of the Roman war against Jerusalem. St Hilary of
Poitiers and St Jerome read them also as words to every hearer, and the rule
they draw from the housetop is that one who has gone up must not go back down
to what he has left. That rule is the subject of this interpretation. Three
of the Mass's prayers ask for what keeping it takes: the Collect for roused
wills, ``arouse \dots\ the wills of Thy faithful''; the Epistle for
perseverance, that the Colossians may walk worthy of God in patience; and the
Secret for release from what draws a man down, that, ``being delivered from
earthly desires, we may pass on to the enjoyments of heaven.'' The English of
the Collect and the Secret is Lasance's.

\subsection{The war that Chrysostom saw}
\label{sec:flight-chrysostom}

Chrysostom holds to the letter, and gives his reasons. The sabbath shows to
whom the Lord is speaking: ``Seest thou that His discourse is addressed to the
Jews \dots\ the apostles surely were not to keep the Sabbath day, neither to
be there, when Vespasian did those things'' (\work{Hom.\ in Mt.}~76.1). The
tribulation without equal is the siege, and anyone who doubts it is to ``study
the writings of Josephus.'' ``No flesh'' means ``no Jewish flesh'': ``If, saith
He, the war of the Romans against the city had prevailed further, all the Jews
had perished.'' The elect, for whose sake the days were shortened, are ``The
believers that were shut up in the midst of them'' (76.2). St Thomas Aquinas
follows the same historical line: the believers left the city and went to the
realm of Agrippa, and the man on the roof and the man in the field are told
to leave everything, since a man will give all he has for his life
(\work{Super Matthaeum}, cap.~24). St Bede, commenting on Mark's account of
the same discourse, records that the Christians of the province, warned by an
oracle, withdrew across the Jordan to the city of Pella, \latin{ut
ecclesiastica narrat Historia}, ``as the Church History relates'' (\work{In
Marci evangelium expositio}~IV, on 13:14).

\subsection{The housetop is the height of the house}
\label{sec:flight-hilary}

Hilary takes the abomination for Antichrist, and reads the flight from Judea
as a flight from the contagion of those who will receive him (\work{In
Matthaeum}~25.3--4). At the housetop he stops. \latin{Praeceptum istud secundum
humanam intelligentiam rationem dicti factique non recipit. Qui enim in tecto
est \dots\ nisi descenderit ex eo abire non poterit}, ``This precept, by human
understanding, does not admit a reason for what is said and done. For whoever
is on the roof cannot leave unless he comes down from it.'' Nor is there any
advantage in standing on the roof rather than staying inside the house. The
words must therefore be weighed \latin{ut momenta praeceptorum coelestium
consequamur}, ``so that we may reach the force of heavenly precepts''
(25.5).

His answer begins from what a roof is: \latin{Tectum est domus fastigium, et
habitationis totius celsa perfectio}, ``The roof is the summit of the house,
and the lofty completion of the whole dwelling.'' The one who stands on it is
the one who stands \latin{in corporis sui perfectione, regeneratione novus,
spiritu celsus}, ``in the perfection of his body, new by rebirth, lofty in
spirit,'' and he \latin{non descendere in humiliora rerum saecularium
cupiditate debebit}, ``ought not to go down into lower things by desire for
worldly things.'' The man in the field, at work on the commandment, is not to
go back to his former cares, nor to take up again the tunic \latin{veterum
peccaminum}, ``of old sins,'' in which he was once clothed.

The woe to women with child cannot, Hilary says, be a curse on a condition of
nature: \latin{absit \dots\ ut quidquam sit, quod homini nisi sua culpa malum
fiat}, ``far be it that anything should be evil to a man except by his own
fault.'' The pregnant are souls heavy with sins, who can neither escape the
storm on the roof nor in the field; the nursing are souls still fed on milk,
who \latin{nec peccata effugerint, nec cibum veri panis acceperint}, ``have
neither escaped their sins nor received the food of the true bread'' (25.6).
Winter and sabbath are \latin{ne aut in peccatorum frigore, aut in otio
bonorum operum reperiamur}, ``that we be not found either in the cold of sins
or in idleness from good works'' (25.7).

\subsection{The perverse doctrine and the eternal mountains}
\label{sec:flight-jerome}

St Jerome lists the plain senses of the abomination and then adds
another: \latin{Abominatio desolationis intelligi potest, et omne dogma
perversum}, ``The abomination of desolation can be understood also as every
perverse doctrine.'' When it stands in the holy place, \latin{hoc est, in
Ecclesia}, ``that is, in the Church,'' we must flee from Judea to the
mountains, \latin{dimissa occidente littera \dots\ appropinquare montibus
aeternis}, ``leaving the letter that kills, to draw near to the eternal
mountains'' (Ps 75:5); we must stay on the roof, where the devil's fiery darts
cannot reach, and not go down \latin{de domo conversationis pristinae}, ``from
the house of our former life'' (\work{In Matthaeum}~IV, on 24:15--18). He
reads the winter and the sabbath first of the captivity of Jerusalem under
Titus and Vespasian, when the cold and the sabbath law would hinder flight
from the city, and then: \latin{Si autem
de consummatione mundi intelligitur, hoc praecipit, ut non refrigescat fides
nostra et in Christum charitas, neque ut otiosi in opere Dei torpeamus
virtutum sabbato}, ``But if it is understood of the consummation of the world,
it commands that our faith and our charity towards Christ not grow cold, and
that we not lie idle and sluggish in the work of God, on a sabbath from the
virtues'' (on 24:20). The days are shortened, he says, \latin{ne temporum mora,
fides concutiatur credentium}, ``lest the faith of believers be shaken by the
length of the time'' (on 24:22).

St Bede, after giving the history of Pella, gives a spiritual sense of his own
that runs close to Jerome's. When heresies and crimes reign among those who
seemed consecrated to the heavenly mysteries, those who stand in Judea, that
is, in the confession of the true faith, \latin{tanto altius virtutum culmen
debemus ascendere, quanto plures ampla et errabunda vitiorum sequi itinera
videmus}, ``ought to climb so much the higher to the summit of the virtues, as
we see more following the broad and wandering ways of the vices.'' The one on
the housetop is \latin{qui excedens animo carnalia, tanquam in aura libera
spiritualiter vivit}, ``he who, rising in mind above carnal things, lives
spiritually as in free air''; he is not to go down to the lowest acts of his
former life, \latin{neque ea quae reliquerat mundi carnisve desideria repetat},
``nor seek again the desires of the world or the flesh which he had left''
(\work{In Marci evangelium expositio}~IV, on 13:15--16). The house he would go
down into is the flesh, the ``earthly house of this habitation'' (2~Cor 5:1).
St Augustine gives the winter and sabbath two senses: any hindrance at all,
\latin{Vel aliud, ne in tristitia aut laetitia rerum temporalium quis inveniatur
a die illa}, ``Or another thing: that no one be found by that day in sorrow or
joy over temporal things'' (\work{Quaestiones Evangeliorum}~I.37).

\subsection{The deceivers in the desert and the inner rooms}
\label{sec:flight-deceivers}

The flight is followed by a second warning: ``If therefore they shall say to
you, Behold he is in the desert: go ye not out. Behold he is in the closets:
believe it not'' (Mt 24:26). Jerome reads the desert as the dogma of the
philosophers and the inner rooms as the hiding places of heretics (on
24:25--26). Augustine makes the lightning the Church itself, spread through
the whole world, \latin{Convenienter enim Ecclesiam nunc fulgur nominavit},
``For he fittingly named the Church lightning here,'' and the ``here'' and
``there'' the schisms that hold some part of the earth, \latin{partem aliquam
tenens}, while the desert and the inner rooms are the \latin{obscuris atque
occultis conventiculis}, ``dark and hidden meetings,'' of heretics
(\work{Quaestiones Evangeliorum}~I.38). The one who has gone up to the
mountains is not to go out after any of them. The letter of Jeremias, whose
words the Introit sings, gave the exiles a like warning: ``Let not your
prophets that are in the midst of you, and your diviners deceive you'' (Jer
29:8).

\subsection{Walking worthy, in patience}
\label{sec:flight-epistle}

The Epistle is a prayer that the Colossians may keep going. Paul asks that
they ``walk worthy of God, in all things pleasing,'' fruitful and growing, and
``Strengthened with all might according to the power of his glory, in all
patience and longsuffering with joy'' (Col 1:10--11). St John Chrysostom notes
that ``with faith he always couples conduct,'' and that ``he that knows
without doing, is even in the way to punishment.'' ``To walk,'' he says, means
``always, not once, but continually,'' and the fruitfulness means being
``always advancing, and nowhere standing still.'' The strength Paul asks for
concerns ``trials and persecutions,'' and the two virtues are distinct:
``long-suffering towards one another, patience towards those without''
(\work{Hom.\ in Col.}~2). St Thomas Aquinas distinguishes them otherwise:
patience is against the hardness of adversity, longsuffering against the delay
of the reward, and ``with joy'' against bearing them in sadness (\work{Super
Colossenses}, cap.~1, lect.~3). He adds that \latin{ad virtutem non sufficit
scire vel velle, nisi immobiliter operetur}, ``for virtue it is not enough to
know or to will, unless one works without wavering.'' Dom Lucien Fromage,
commenting on this lesson, writes that ``The just man cannot possibly remain
stationary in this world; he must either descend or ascend'' (\work{The
Liturgical Year}, vol.~11, p.~505, citing Ps 83:6).

The Gradual's deliverance is, for Aquinas, partly present already. On
\latin{salvasti \dots\ ex affligentibus nos} he writes, \latin{vel in praesenti:
quia etsi corporaliter nobis noceant, spiritualiter tamen nocere non possunt},
``either in the present, because though they may harm us in body, yet in spirit
they cannot harm us,'' or in the future (\work{Super Psalmos}~43, nn.~3--4).
The devil's fiery darts are aimed at the spirit, and the one who stays on
Jerome's roof, out of their reach, has that present deliverance already.

\subsection{The prayer of those who will not come down}
\label{sec:flight-orations}

The Collect asks God to ``arouse \dots\ the wills of Thy faithful,'' that,
``by more earnestly following after the fruit of the divine work, they may
the more abundantly partake of Thy mercies'' (Lasance). Jerome's danger was
that we \latin{otiosi in opere Dei torpeamus}, ``lie idle and sluggish in the
work of God.'' The Collect uses the same noun, \latin{divíni óperis fructum},
``the fruit of the divine work,'' and asks for the opposite of idleness:
wills roused, and the fruit followed \latin{propénsius}, ``more readily,''
where Lasance has ``more earnestly.'' A comparative asks for more than
staying on the roof.

The Secret asks that, ``being delivered from earthly desires, we may pass on
to the enjoyments of heaven'' (Lasance). Hilary's man on the housetop must
not go down \latin{rerum saecularium cupiditate}, ``by desire for worldly
things,'' and Bede's must not seek again the \latin{desideria} of the world
and the flesh that he has left. Both commands assume that the one who keeps
them has already left his former desires. The Secret asks that leaving for
\latin{ómnium nostrum}, ``all of us,'' and its Latin has the two nouns that
Hilary and Bede use: delivered from earthly \latin{cupiditátibus}, we are to
pass over to heavenly \latin{desidéria}. Its verb, \latin{transeámus}, ``may
we pass over,'' is a verb of travel, and St Gregory the Great, preaching on
Luke's account of the same discourse, asks, \latin{Quid est vita mortalis,
nisi via?}, ``What is mortal life but a way?'', and wonders at those who grow
weary on the way and yet do not wish it to end
(\work{Homiliae in Evangelia}~I.1.3). His counsel at the end of the sermon is
\latin{Nolite ergo, fratres mei, diligere quem videtis diu stare non posse},
``Do not then, my brethren, love what you see cannot stand for long'' (I.1.5,
citing 1~Jn 2:15).

The Gospel's only command to pray is ``pray that your flight be not in the
winter or on the sabbath.'' Aquinas comments that the Lord bids them pray
\latin{quia in talibus ad Deum solum est recurrendum}, ``because in such
things one must have recourse to God alone.'' The Alleluia and the
Offertory, ``From the depths I have cried,'' are prayer of that kind, a cry
that Hilary pictures made from the surge of human calamities. Mark gives the
Communion's promise, ``whatsoever you ask when you pray, believe that you
shall receive,'' in another setting, after the withered fig tree (Mk
11:20--24). Its ``whatsoever'' excepts no prayer, and so it reaches the
prayer the Gospel commands of those who flee.

Those who say the Postcommunion have received the food of the true bread
that Hilary's nursing souls lack. What they ask in it, that ``whatever is evil
in our hearts may be restored by its gift of healing'' (Lasance), is the other
thing those souls lack, escape from their sins.

\subsection{The letter and the hearer}
\label{sec:flight-difficulty}

The difficulty is one of method, and the Fathers do not settle it. Hilary
argues that the housetop precept, taken literally, has no reason, and goes
from the letter to the soul. Chrysostom does not leave the letter: the words
were spoken to Jews who would see the city taken, and the apostles would not
be there. Both cannot be the whole account. Chrysostom's reading would make
the passage a fulfilled history, to be read on this Sunday as proof that what
the Lord foretold came true. Hilary's and Jerome's make it a command spoken to
this congregation, who are on the roof by baptism and in the field at their
work, and who must not come down or go back.

This interpretation follows Hilary and Jerome, and it does so with
Chrysostom's sense standing first. Jerome himself sets out the plain senses
before he gives the spiritual one, and Bede gives the flight to Pella
before he gives the climb to the summit of the virtues. The history shows
that the command was obeyed and that those who obeyed it were saved; the
spiritual sense applies the same command to those who hear it now. A second
difference remains. Hilary refers the abomination to Antichrist; Jerome
leaves several referents open and adds every perverse doctrine.

\begin{fourSenses}
\item[Literal.] Jesus warns that when the abomination Daniel spoke of stands
in the holy place, those in Judea must flee at once to the mountains, not
turning back for anything, and pray that the flight not fall in winter or on
the sabbath (Mt 24:15--20); Chrysostom and Aquinas read it of the Roman war,
Bede of the flight to Pella. Paul prays that the Colossians walk worthy of
God in patience and longsuffering (Col 1:10--11).
\item[Allegorical.] The abomination is Antichrist (Hilary) or every perverse
doctrine standing in the Church (Jerome); Judea is the confession of the true
faith (Bede); the mountains are the eternal mountains (Jerome); the roof is
the height of the house, the perfection of one reborn (Hilary); the house is
the flesh (Bede); the lightning is the Church spread through the world, and
the desert and the inner rooms the hidden meetings of heretics (Augustine).
\item[Moral.] Not to go down from the roof by desire for worldly things, nor
go back for the tunic of old sins (Hilary); not to let faith and charity grow
cold nor lie idle in the work of God (Jerome); not to be found by that day in
sorrow or joy over temporal things (Augustine); to climb the higher the more
others wander (Bede); to walk always and never stand still (Chrysostom).
\item[Anagogical.] The days are shortened for the elect, so that the faith of
believers is not shaken by the length of the time (Jerome); and on that day
the one who did not go down is found neither in the cold of sins nor idle from
good works (Hilary, 25.7).
\end{fourSenses}
